\subsubsection{Die Ausgangslage – dotCloud und das „Works on my Machine"-Problem}

Bevor Docker existierte, stand die Softwareentwicklung vor einem alltäglichen 
Problem: Eine Anwendung funktionierte auf dem Rechner des Entwicklers 
einwandfrei, versagte jedoch in der Test- oder Produktionsumgebung. 
Unterschiedliche Betriebssystemversionen, abweichende Bibliotheken oder 
fehlende Abhängigkeiten führten zu schwer nachvollziehbaren Fehlern. 
Dieses Phänomen war so verbreitet, dass es in der Entwickler-Community 
seinen eigenen Begriff bekam: \textit{„Works on my machine"} 
\cite{ionos_docker}.

Solomon Hykes, ein französisch-amerikanischer Softwareentwickler, 
gründete 2008 gemeinsam mit Sebastien Pahl das Unternehmen 
\textbf{dotCloud} in San Francisco \cite{datacenter_docker}. dotCloud 
war ein Platform-as-a-Service (PaaS)-Anbieter, der Entwicklern ermöglichte, 
ihre Anwendungen auf der dotCloud-Infrastruktur zu deployen, ohne sich 
selbst um Server kümmern zu müssen. Das klingt einfach, war es aber 
nicht: dotCloud musste Anwendungen von hunderten verschiedenen Kunden 
gleichzeitig betreiben, jede mit eigenen Abhängigkeiten, Laufzeitumgebungen 
und Konfigurationen, und das alles isoliert voneinander auf denselben 
physischen Servern \cite{datacenter_docker}.

Um dieses Problem intern zu lösen, begann das dotCloud-Team mit der 
Entwicklung einer Technologie, die Anwendungen zusammen mit ihren 
gesamten Abhängigkeiten in portable, isolierte Einheiten verpackte. 
Diese Einheiten, später als \textbf{Container} bekannt, konnten auf 
jeder Umgebung gleich ausgeführt werden, unabhängig vom darunterliegenden 
System \cite{ionos_docker}.

\subsubsection{Die technische Basis – Linux Namespaces und cgroups}

Docker war keine vollständig neue Erfindung, sondern baute auf bereits 
bestehenden Linux-Kernelfunktionen auf, die bis dahin jedoch kaum 
zugänglich oder bekannt waren \cite{ionos_docker}.

Die zwei wichtigsten Bausteine waren:

\begin{itemize}
    \item \textbf{Linux Namespaces}: Namespaces isolieren Prozesse 
    voneinander. Jeder Container bekommt seine eigene Sicht auf das 
    System, also eigene Prozess-IDs, eigenes Netzwerk und eigenes 
    Dateisystem. Aus Sicht des Containers läuft er auf einem vollständig 
    isolierten System, obwohl er sich den Kernel mit dem Host teilt 
    \cite{ionos_docker}.

    \item \textbf{Control Groups (cgroups)}: cgroups wurden von Google 
    entwickelt und 2008 in den Linux-Kernel integriert. Sie ermöglichen 
    es, den Ressourcenverbrauch von Prozessen zu begrenzen und zu 
    überwachen, also wie viel CPU, Arbeitsspeicher oder Netzwerkbandbreite 
    ein Container maximal nutzen darf \cite{ionos_docker}.
\end{itemize}

Diese Kombination aus Namespaces und cgroups bildete das Fundament, 
auf dem Docker aufbaute. Der entscheidende Beitrag von Docker war 
nicht die Erfindung dieser Technologien, sondern ihre \textbf{Zugänglichkeit}: 
Docker verpackte die Komplexität dieser Linux-Mechanismen in eine 
einfache, entwicklerfreundliche Oberfläche \cite{datacenter_docker}.

\subsubsection{Die PyCon-Demo 2013 – ein Wendepunkt}

Am 21. März 2013 präsentierte Solomon Hykes Docker zum ersten Mal 
öffentlich auf der \textbf{PyCon-Konferenz} in Santa Clara, Kalifornien 
\cite{datacenter_docker}. Die Demo dauerte nur wenige Minuten, Hykes 
zeigte live, wie sich eine Anwendung in einem Container starten, stoppen 
und auf einem anderen System ausführen ließ. Die Reaktion des Publikums 
war unmittelbar begeistert.

Noch am selben Tag wurde Docker als \textbf{Open-Source-Projekt} auf 
GitHub veröffentlicht. Innerhalb weniger Stunden verbreitete sich die 
Demo viral in der Entwickler-Community. Was als internes Werkzeug von 
dotCloud begonnen hatte, wurde über Nacht zu einem der meistdiskutierten 
Projekte in der Softwareentwicklung \cite{datacenter_docker}.

\subsubsection{Von dotCloud zu Docker Inc. – 2013 bis 2015}

Der Erfolg der PyCon-Demo hatte unmittelbare Konsequenzen für dotCloud. 
Das Interesse an Docker überwältigte das ursprüngliche PaaS-Geschäft 
bei weitem. Innerhalb weniger Monate zählte das GitHub-Repository 
tausende Sterne und hunderte externe Mitwirkende \cite{datacenter_docker}.

Im Oktober 2013 traf Hykes eine strategische Entscheidung: dotCloud 
wurde in \textbf{Docker Inc.} umbenannt, und das Unternehmen richtete 
seinen gesamten Fokus auf die Weiterentwicklung von Docker. Das 
ursprüngliche PaaS-Geschäft von dotCloud wurde später verkauft 
\cite{datacenter_docker}.

2014 und 2015 folgten weitere wichtige Meilensteine:

\begin{itemize}
    \item \textbf{Docker Hub} wurde eingeführt, eine öffentliche 
    Registry, in der Entwickler fertige Container-Images teilen und 
    wiederverwenden konnten \cite{ionos_docker}.
    
    \item Große Technologieunternehmen wie \textbf{Google, Microsoft 
    und Amazon} kündigten offizielle Docker-Unterstützung in ihren 
    Cloud-Plattformen an \cite{ionos_docker}.
    
    \item Docker erhielt massive Investitionen, unter anderem eine 
    Finanzierungsrunde von 40 Millionen US-Dollar im Jahr 2014 
    \cite{datacenter_docker}.
\end{itemize}

\subsubsection{Die Open Container Initiative – ein offener Standard}

Mit dem rasanten Wachstum von Docker entstand auch die Frage nach 
Standardisierung. 2015 gründeten Docker Inc. gemeinsam mit Unternehmen 
wie Google, Microsoft, Red Hat und anderen die 
\textbf{Open Container Initiative (OCI)} unter dem Dach der Linux 
Foundation \cite{datacenter_docker}.

Das Ziel der OCI war es, offene, herstellerunabhängige Standards für 
Container-Formate und Laufzeitumgebungen zu definieren. Damit sollte 
sichergestellt werden, dass Container nicht an eine bestimmte Plattform 
oder einen bestimmten Anbieter gebunden sind \cite{ionos_docker}. 
Die OCI veröffentlichte zwei zentrale Spezifikationen:

\begin{itemize}
    \item \textbf{image-spec}: definiert das Format eines 
    Container-Images
    \item \textbf{runtime-spec}: definiert wie ein Container aus 
    einem Image gestartet wird
\end{itemize}

\subsubsection{Docker heute}

Heute ist Docker aus der modernen Softwareentwicklung nicht mehr 
wegzudenken. Laut einer Umfrage von Stack Overflow aus dem Jahr 2023 
gehört Docker zu den meistgenutzten Werkzeugen unter Entwicklern 
weltweit \cite{ionos_docker}.

Docker Inc. selbst durchlief in den letzten Jahren mehrere 
Veränderungen, darunter finanzielle Schwierigkeiten und eine 
Neuausrichtung des Unternehmens. Dennoch blieb die Docker-Technologie 
selbst stabil und wird aktiv weiterentwickelt. Die Container-Runtime 
\textbf{containerd}, die ursprünglich Teil von Docker war, wurde 2017 
als eigenständiges Open-Source-Projekt ausgegliedert und ist heute 
die Basis vieler moderner Container-Plattformen, darunter auch 
Kubernetes \cite{ionos_docker}.

Die Grundidee von 2013 ist dabei dieselbe geblieben: Anwendungen 
portabel, reproduzierbar und isoliert bereitzustellen, unabhängig 
davon, wo sie laufen.